% !TEX program = xelatex
% Arabic Beamer presentation with RTL slides using polyglossia.
% Title slide + 2 content slides. The bidi package handles
% right-to-left layout automatically within Beamer frames.
\documentclass[aspectratio=169]{beamer}

\usepackage{fontspec}
\usepackage{polyglossia}
\setmainlanguage[locale=mashriq]{arabic}
\setotherlanguage{english}

\usepackage{bidi}

\newfontfamily\arabicfont[Script=Arabic]{Amiri}
\setbeamertemplate{navigation symbols}{}
\setbeamercolor{frametitle}{bg=blue!20, fg=black}
\setbeamercolor{title}{fg=blue!70!black}

\title{التنضيدُ العربيّ في \LR{Beamer}}
\author{د. نبراس أبو الذهب}
\institute{مركز نبراس}
\date{2025}

\begin{document}

\begin{frame}
  \titlepage
\end{frame}

\begin{frame}{مقدمة}
  يَدعمُ \LR{Beamer} اللغةَ العربيةَ عبرَ حزمتي:
  \begin{itemize}
    \item \LR{polyglossia} لإعدادِ اللغة.
    \item \LR{bidi} لإدارةِ الاتجاهِ من اليمينِ إلى اليسار.
  \end{itemize}
  تَتولّى \LR{bidi} عكسَ اتجاهِ النصوصِ والقوائمِ تلقائياً.
\end{frame}

\begin{frame}{النقاطُ الرئيسية}
  \begin{itemize}
    \item استخدمْ \LR{XeLaTeX} أو \LR{LuaLaTeX} لتشغيلِ \LR{Beamer} العربيّ.
    \item عيّنْ خطّاً عربيّاً عبرَ \LR{\textbackslash newfontfamily}.
    \item استخدمْ \LR{\textbackslash LR\{\}} لإدراجِ نصوصٍ إنجليزية.
    \item يَعملُ اتجاهُ الشرائحِ من اليمينِ إلى اليسارِ تلقائياً.
  \end{itemize}
\end{frame}

\end{document}
